% GENERATED FILE. Edit canonical lessons, then run npm run generate:books.
% canonical-source-hash: fnv1a64:111ec894d0707b6a
% canonical-lessons: MR-C442-sikhar, MR-C442-dari, MR-C442-kada, MR-C442-paytha, MR-C442-dongarrang, MR-R442-review-436-442

\chapter{Peaks and Valleys}
\label{ch:mr-peaks-and-valleys}
\hlchaptermodality{442}

\begin{chapteropening}
\textbf{By the end of this chapter:} \emph{I can talk about peaks, valleys, cliffs, the foot of a hill and a mountain range.}

Complete the last lesson of chapter 442: I can talk about peaks, valleys, cliffs, the foot of a hill and a mountain range.
\end{chapteropening}

\section[\texorpdfstring{\mr{शिखर} (śikhar)}{शिखर (śikhar)}]{\mr{शिखर} (śikhar) — a peak; a summit}

\label{lesson:MR-C442-sikhar}



\begin{quote}
Before the new one: say the Marathi for to set, then the Marathi for to shine.
\end{quote}

\subsection*{You'll want to know: \mr{शिखर}}
\textbf{\mr{शिखर}} — \emph{śikhar} — \textquotedblleft{}a peak; a summit\textquotedblright{}.

\begin{quote}
\mr{आम्ही} \mr{पहाटे} \mr{शिखरावर} \mr{पोहोचलो}.
\end{quote}

\begin{quote}
\emph{āmhī pahāṭe śikharāvar pohoclo.}
\end{quote}

\begin{quote}
We reached the summit at dawn.
\end{quote}

Neuter.

Also the top of a temple, and the height of success: \mr{यशाचं} \mr{शिखर}.

\begin{cousinweb}
\textbf{\mr{शिखर}} (\emph{śikhar}) is a tatsama, taken from Sanskrit \emph{śikhara}, \textquotedblleft{}a peak, a summit, a pinnacle\textquotedblright{} (Monier-Williams).
\end{cousinweb}

\begin{grammarlens}[title={What you've built}]
One more word for getting around.
\end{grammarlens}

\subsection*{Your turn}
\begin{itemize}
\raggedright
  \item \emph{Say it:} \emph{śikhar}
  \item \emph{Say it:} \emph{śikhar}, once more
  \item \emph{Say it:} \emph{āmhī pahāṭe śikharāvar pohoclo.} — We reached the summit at dawn.
  \item \emph{Recall:} say the Marathi for to set, then the Marathi for to shine, then say \emph{śikhar} again
\end{itemize}

\begin{tcolorbox}[breakable,colback=teal!4,colframe=teal!35!black,title={Before you move on}]
What does \mr{शिखर} mean? (\textquotedblleft{}A peak; a summit\textquotedblright{}.) Say it once more.
\end{tcolorbox}
\section[\texorpdfstring{\mr{दरी} (darī)}{दरी (darī)}]{\mr{दरी} (darī) — a valley; a gorge}

\label{lesson:MR-C442-dari}



\begin{quote}
Before the new one: say the Marathi for to shine, then the Marathi for a peak.
\end{quote}

\subsection*{You'll want to know: \mr{दरी}}
\textbf{\mr{दरी}} — \emph{darī} — \textquotedblleft{}a valley; a gorge\textquotedblright{}.

\begin{quote}
\mr{शिखरावरून} \mr{खाली} \mr{खोल} \mr{दरी} \mr{दिसते}.
\end{quote}

\begin{quote}
\emph{śikharāvarūn khālī khol darī diste.}
\end{quote}

\begin{quote}
From the summit you can see a deep valley below.
\end{quote}

Feminine.

\mr{खोल} \mr{दरी}, a deep valley.

\begin{cousinweb}
\textbf{\mr{दरी}} (\emph{darī}) is inherited: CDIAL lists Marathi \emph{darī}, \textquotedblleft{}chasm\textquotedblright{}, under Sanskrit \emph{dara}, \emph{darī}, \textquotedblleft{}a cleft, a cave\textquotedblright{} (CDIAL 6188).
\end{cousinweb}

\begin{grammarlens}[title={What you've built}]
One more word for getting around.
\end{grammarlens}

\subsection*{Your turn}
\begin{itemize}
\raggedright
  \item \emph{Say it:} \emph{darī}
  \item \emph{Say it:} \emph{darī}, once more
  \item \emph{Say it:} \emph{śikharāvarūn khālī khol darī diste.} — From the summit you can see a deep valley below.
  \item \emph{Recall:} say the Marathi for to shine, then the Marathi for a peak, then say \emph{darī} again
\end{itemize}

\begin{tcolorbox}[breakable,colback=teal!4,colframe=teal!35!black,title={Before you move on}]
What does \mr{दरी} mean? (\textquotedblleft{}A valley; a gorge\textquotedblright{}.) Say it once more.
\end{tcolorbox}
\section[\texorpdfstring{\mr{कडा} (kaḍā)}{कडा (kaḍā)}]{\mr{कडा} (kaḍā) — a cliff; an edge}

\label{lesson:MR-C442-kada}



\begin{quote}
Before the new one: say the Marathi for a peak, then the Marathi for a valley.
\end{quote}

\subsection*{You'll want to know: \mr{कडा}}
\textbf{\mr{कडा}} — \emph{kaḍā} — \textquotedblleft{}a cliff; an edge\textquotedblright{}.

\begin{quote}
\mr{कड्यावरून} \mr{धबधबा} \mr{खाली} \mr{पडतो}.
\end{quote}

\begin{quote}
\emph{kaḍyāvarūn dhabdhabā khālī paḍto.}
\end{quote}

\begin{quote}
The waterfall drops down from the cliff.
\end{quote}

Masculine.

Before an ending \mr{कडा} becomes \mr{कड्या}-: \mr{कड्यावरून}, from the cliff.

\begin{cousinweb}
\textbf{\mr{कडा}} (\emph{kaḍā}) is inherited: CDIAL lists Marathi \emph{kaḍā}, \textquotedblleft{}precipice\textquotedblright{}, under Sanskrit \emph{kaḍāra} (CDIAL 2655).
\end{cousinweb}

\begin{grammarlens}[title={What you've built}]
One more word for getting around.
\end{grammarlens}

\subsection*{Your turn}
\begin{itemize}
\raggedright
  \item \emph{Say it:} \emph{kaḍā}
  \item \emph{Say it:} \emph{kaḍā}, once more
  \item \emph{Say it:} \emph{kaḍyāvarūn dhabdhabā khālī paḍto.} — The waterfall drops down from the cliff.
  \item \emph{Recall:} say the Marathi for a peak, then the Marathi for a valley, then say \emph{kaḍā} again
\end{itemize}

\begin{tcolorbox}[breakable,colback=teal!4,colframe=teal!35!black,title={Before you move on}]
What does \mr{कडा} mean? (\textquotedblleft{}A cliff; an edge\textquotedblright{}.) Say it once more.
\end{tcolorbox}
\section[\texorpdfstring{\mr{पायथा} (pāythā)}{पायथा (pāythā)}]{\mr{पायथा} (pāythā) — a foothill; the foot of a hill}

\label{lesson:MR-C442-paytha}



\begin{quote}
Before the new one: say the Marathi for a valley, then the Marathi for a cliff.
\end{quote}

\subsection*{You'll want to know: \mr{पायथा}}
\textbf{\mr{पायथा}} — \emph{pāythā} — \textquotedblleft{}a foothill; the foot of a hill\textquotedblright{}.

\begin{quote}
\mr{डोंगराच्या} \mr{पायथ्याशी} \mr{एक} \mr{लहान} \mr{गाव} \mr{आहे}.
\end{quote}

\begin{quote}
\emph{ḍoṅgarācyā pāyathyāśī ek lahān gāv āhe.}
\end{quote}

\begin{quote}
There is a small village at the foot of the hill.
\end{quote}

Masculine.

\mr{पायथ्याशी}, at the foot; the opposite is the \mr{शिखर}.

\begin{cousinweb}
\textbf{\mr{पायथा}} (\emph{pāythā}) is inherited: CDIAL lists Marathi \emph{pāythā}, \textquotedblleft{}extremity of the feet\textquotedblright{}, under Sanskrit \emph{pādānta}, \textquotedblleft{}the end at the feet\textquotedblright{} (CDIAL 8072).
\end{cousinweb}

\begin{grammarlens}[title={What you've built}]
One more word for getting around.
\end{grammarlens}

\subsection*{Your turn}
\begin{itemize}
\raggedright
  \item \emph{Say it:} \emph{pāythā}
  \item \emph{Say it:} \emph{pāythā}, once more
  \item \emph{Say it:} \emph{ḍoṅgarācyā pāyathyāśī ek lahān gāv āhe.} — There is a small village at the foot of the hill.
  \item \emph{Recall:} say the Marathi for a valley, then the Marathi for a cliff, then say \emph{pāythā} again
\end{itemize}

\begin{tcolorbox}[breakable,colback=teal!4,colframe=teal!35!black,title={Before you move on}]
What does \mr{पायथा} mean? (\textquotedblleft{}A foothill; the foot of a hill\textquotedblright{}.) Say it once more.
\end{tcolorbox}
\section[\texorpdfstring{\mr{डोंगररांग} (ḍoṅgarrāṅg)}{डोंगररांग (ḍoṅgarrāṅg)}]{\mr{डोंगररांग} (ḍoṅgarrāṅg) — a mountain range}

\label{lesson:MR-C442-dongarrang}



\begin{quote}
Before the new one: say the Marathi for a cliff, then the Marathi for a foothill.
\end{quote}

\subsection*{You'll want to know: \mr{डोंगररांग}}
\textbf{\mr{डोंगररांग}} — \emph{ḍoṅgarrāṅg} — \textquotedblleft{}a mountain range\textquotedblright{}.

\begin{quote}
\mr{पश्चिमेला} \mr{एक} \mr{मोठी} \mr{डोंगररांग} \mr{आहे}.
\end{quote}

\begin{quote}
\emph{paścimelā ek moṭhī ḍoṅgarrāṅg āhe.}
\end{quote}

\begin{quote}
There is a great mountain range to the west.
\end{quote}

Feminine.

\mr{डोंगर}, a hill, and \mr{रांग}, a row.

\begin{cousinweb}
\textbf{\mr{डोंगररांग}} (\emph{ḍoṅgarrāṅg}) is built from \textbf{\mr{डोंगर}}, \textquotedblleft{}a hill\textquotedblright{}, and \textbf{\mr{रांग}}, \textquotedblleft{}a row, a line\textquotedblright{}, both of uncertain origin.
\end{cousinweb}

\begin{grammarlens}[title={What you've built}]
One more word for getting around.
\end{grammarlens}

\subsection*{Your turn}
\begin{itemize}
\raggedright
  \item \emph{Say it:} \emph{ḍoṅgarrāṅg}
  \item \emph{Say it:} \emph{ḍoṅgarrāṅg}, once more
  \item \emph{Say it:} \emph{paścimelā ek moṭhī ḍoṅgarrāṅg āhe.} — There is a great mountain range to the west.
  \item \emph{Recall:} say the Marathi for a cliff, then the Marathi for a foothill, then say \emph{ḍoṅgarrāṅg} again
\end{itemize}

\begin{tcolorbox}[breakable,colback=teal!4,colframe=teal!35!black,title={Before you move on}]
What does \mr{डोंगररांग} mean? (\textquotedblleft{}A mountain range\textquotedblright{}.) Say it once more.
\end{tcolorbox}
\section[(dialogue)]{Looking back: from Heat and Humidity to Peaks and Valleys}

\label{lesson:MR-R442-review-436-442}



\begin{quote}
Say the Marathi for the atmosphere and great heat.
\end{quote}

\subsection*{Your turn}
Say these aloud:

\begin{itemize}
\raggedright
  \item the atmosphere, great heat, sultry heat, humid, to blaze
  \item coolness, a snowfall, to sting, to freeze, to melt
  \item a drizzle, torrential, a pre-monsoon storm, to pour down, to get wet
  \item a cyclone, a whirlwind, a rumble of thunder, to howl, to crackle
  \item a drought, a cloudburst, a severe spell, to lash, to recede
  \item the horizon, moonlight, to rise, to set, to shine
  \item a peak, a valley, a cliff, a foothill, a mountain range
\end{itemize}

\begin{tcolorbox}[breakable,colback=teal!4,colframe=teal!35!black,title={Before you move on}]
The atmosphere? (\textbf{\mr{वातावरण}}, \emph{vātāvaraṇ}.) Great heat? (\textbf{\mr{उष्णता}}, \emph{uṣṇatā}.) Sultry heat? (\textbf{\mr{उकाडा}}, \emph{ukāḍā}.) Humid? (\textbf{\mr{दमट}}, \emph{damaṭ}.) To blaze? (\textbf{\mr{तळपणे}}, \emph{taḷapṇe}.) Coolness? (\textbf{\mr{गारवा}}, \emph{gārvā}.) A snowfall? (\textbf{\mr{हिमवर्षाव}}, \emph{himvarṣāv}.) To sting? (\textbf{\mr{झोंबणे}}, \emph{jhombṇe}.) To freeze? (\textbf{\mr{गोठणे}}, \emph{goṭhṇe}.) To melt? (\textbf{\mr{वितळणे}}, \emph{vitaḷṇe}.) A drizzle? (\textbf{\mr{रिमझिम}}, \emph{rimjhim}.) Torrential? (\textbf{\mr{मुसळधार}}, \emph{musaḷdhār}.) A pre-monsoon storm? (\textbf{\mr{वळीव}}, \emph{vaḷīv}.) To pour down? (\textbf{\mr{बरसणे}}, \emph{barasṇe}.) To get wet? (\textbf{\mr{भिजणे}}, \emph{bhijṇe}.) A cyclone? (\textbf{\mr{चक्रीवादळ}}, \emph{cakrīvādaḷ}.) A whirlwind? (\textbf{\mr{वावटळ}}, \emph{vāvṭaḷ}.) A rumble of thunder? (\textbf{\mr{गडगडाट}}, \emph{gaḍgaḍāṭ}.) To howl? (\textbf{\mr{घोंगावणे}}, \emph{ghoṅgāvṇe}.) To crackle? (\textbf{\mr{कडाडणे}}, \emph{kaḍāḍṇe}.) A drought? (\textbf{\mr{दुष्काळ}}, \emph{duṣkāḷ}.) A cloudburst? (\textbf{\mr{ढगफुटी}}, \emph{ḍhagphuṭī}.) A severe spell? (\textbf{\mr{कडाका}}, \emph{kaḍākā}.) To lash? (\textbf{\mr{झोडपणे}}, \emph{jhoḍapṇe}.) To recede? (\textbf{\mr{ओसरणे}}, \emph{osarṇe}.) The horizon? (\textbf{\mr{क्षितिज}}, \emph{kṣitij}.) Moonlight? (\textbf{\mr{चांदणे}}, \emph{cāndṇe}.) To rise? (\textbf{\mr{उगवणे}}, \emph{ugavṇe}.) To set? (\textbf{\mr{मावळणे}}, \emph{māvaḷṇe}.) To shine? (\textbf{\mr{चमकणे}}, \emph{camakṇe}.) A peak? (\textbf{\mr{शिखर}}, \emph{śikhar}.) A valley? (\textbf{\mr{दरी}}, \emph{darī}.) A cliff? (\textbf{\mr{कडा}}, \emph{kaḍā}.) A foothill? (\textbf{\mr{पायथा}}, \emph{pāythā}.) A mountain range? (\textbf{\mr{डोंगररांग}}, \emph{ḍoṅgarrāṅg}.)
\end{tcolorbox}
